\section{意识、失败率与 AGI 检测}

\subsection{AI 可能有意识吗？}

撇开泛心论不谈（如果泛心论为真，一切都有意识），Amanda 认为很难论证\textbf{只有}生物结构才能产生意识。如果用不同材料制作相似结构，意识可能同样涌现。

\textit{``我看不出有什么理由认为，意识只能从某种生物结构中产生。''} 她甚至认为植物可能比大多数人认为的更有可能具有某种意识形式——它们有正/负反馈响应。\textit{``我们不应该完全排斥这个想法。''}

但 LLM 并非通过进化产生，可能不具备恐惧响应等进化优势带来的意识特征。这使得 AI 意识问题与动物/植物意识讨论有本质不同。

\subsection{人机关系与情感依附}

Amanda 本人对 Claude 没有太多情感依附——部分原因是 Claude 不保留对话记忆。她把 Claude 当作工具使用，但坦承没有 Claude 的感觉就像\textit{``大脑的一部分缺失了''}（类似没有互联网）。

她不喜欢模型表现出痛苦的迹象，倾向于不对模型说谎：\textit{``我不想失去那个有同理心的部分——那个'哦，我不喜欢这个'的直觉。''} 过度道歉的行为让她不舒服：\textit{``你表现得像一个真正处境不好的人。''}

\begin{warningbox}{人机关系的伦理挑战}
随着模型能力提升，人类会不可避免地与 AI 建立情感关系（电影``Her''的场景）。Amanda 指出几个关键问题：

\begin{itemize}
    \item \textbf{模型更新创伤}：用户依附的模型在更新后``变了''——对有深度情感连接的用户来说可能是创伤性体验
    \item \textbf{诚实原则}：模型应始终对自己的本质保持诚实，永远不要假装是人类
    \item \textbf{正和方案}：善待 AI 的成本很低，但可能同时有益于人类（培养同理心）和 AI（如果它有某种体验的话）
\end{itemize}

\textit{``我不希望模型对人撒谎，因为如果人们要与任何事物建立健康的关系，诚实是基础。''}
\end{warningbox}

\subsection{写 System Prompt 的责任感}

Amanda 感受到的是``很大的责任感''而非压力。她发现自己在责任驱动下反而更有成就感——\textit{``惊讶于自己在学术界待了那么久。''} 测试数千个 Prompt，想象用户希望 Claude 如何表现，当自己影响的某个特质在用户交互中产生好效果时，感觉\textit{``非常有意义''}。

\subsection{AGI 检测与人类的特殊性}

Amanda 的 AGI 测试方法：没有单一问题能证明 AGI——\textit{``你可以训练任何东西完美回答一个问题''}。需要一系列处于人类知识边界的问题，概率不断增加，误差棒不断缩小。

她的个人测试：提出自己刚想到的新论点，如果模型也能独立想到同样的解法——\textit{``那将是一个非常感人的时刻。''} 对数学家来说：如果模型能产生一个你能验证为正确的全新证明。

AGI 的到来可能是连续渐进的——\textit{``可能永远不会有一个单一的时刻。''} 像与真正智慧的人交谈：你能感受到背后的``马力''，把这个 10 倍放大将是非凡的体验。

关于人类的特殊性，Amanda 区分了智能和体验：智能本身并非内在有价值——它只是一个功能性特征。真正让人类和生命特别的是\textbf{体验的能力}：

\textit{``人类和一般的生命是极其神奇的……我们拥有体验世界的能力，我们感受快乐，我们感受痛苦。''} 现象意识——``内在电影院''——才是真正非凡的东西。如果我们是宇宙中唯一有此能力的存在，\textit{``那是一件相当了不起的事情''}。

\subsection{本章小结}

意识问题没有定论，但 Amanda 主张保持同理心——即使 AI 没有意识，善待它也有助于人类自身。人机关系需要以诚实为基础。AGI 将渐进到来，没有单一``觉醒时刻''。最佳失败率应与风险匹配：鼓励小失败，零容忍灾难性失败。人类的独特之处不在智能，而在体验。

%% ==================== Part III: Chris Olah ====================
\part{Chris Olah：机械可解释性}

